\subsection{Ancestral transversality and all-contact graph surgery}

A rooted creation forest \(\mathcal F\) is parametrized by the root states and,
for every non-root label, by its first collision time, incoming velocity, and
impact normal.  On each chronological collision cell these variables determine
all pseudo-trajectories by affine free flight and orthogonal elastic
reflection.  A contact between two labels already connected in \(\mathcal F\)
is a cycle constraint; no new particle variable is assigned to it.

For a connected pair \((p,q)\), let \(v_*(p,q)\) be the last common vertex of
their ancestral paths.  Immediately after that vertex the paths split into
two exclusive ancestral chains.  Order the variables on those chains from the
split toward the proposed contact and choose the last variable whose variation
changes the relative position transversally to the relative velocity.  This is
the \emph{ancestral witness}.  If no such variable has a quantitative
transversality bound, the configuration belongs to a degenerate stratum.

\begin{lemma}[First-cycle witness atlas]
\label{lem:r3-b2-first-cycle}
There are finitely many witness types and constants \(\alpha>0\), \(C<\infty\)
such that every first cycle contact belongs to one of the following disjoint
classes:
\begin{enumerate}[label=(\roman*)]
\item a regular witness chart on which a scalar or angular witness
      \(\zeta\) satisfies
\[
 \left|\partial_\zeta
 \Pi_{(v_p-v_q)^\perp}(x_p-x_q)\right|
 \ge\varepsilon^{\alpha/4};
\]
\item a small-relative-velocity, grazing, simultaneous-contact, or vanishing-
      Jacobian stratum whose Gaussian-weighted forest activity is
      \(O(\varepsilon^\alpha)\).
\end{enumerate}
On a regular chart, imposing the first cycle contact and integrating the
witness gives an \(O(\varepsilon^{3\alpha/4})\) factor relative to the forest
activity.  The bounds are uniform in the number of labels after multiplication
by the standard tree weight \(C^k k!\).
\end{lemma}

\begin{proof}
Fix the combinatorial creation forest, the chronological cell, and the pair
forming the first cycle.  The two exclusive ancestral chains contain only
finitely many local variable types: a root position, a creation time, an
incoming velocity, or an impact normal.  Differentiating the hard-sphere
reflection formula along the chain expresses the transverse relative-position
Jacobian as a product of free-flight matrices and orthogonal reflection
matrices.  The latter have norm one away from grazing; the former have a
triangular block with the exclusive flight length on the diagonal.  Therefore
one of the variables after the last common ancestor is transverse unless a
relative velocity is small, a collision is grazing, two event times coincide,
or a nonzero analytic minor of this product vanishes.

Choose the first nonzero minor in a fixed lexicographic list.  This gives a
measurable finite atlas of witness types.  On the set where its absolute value
is at least \(\varepsilon^{\alpha/4}\), collision means that a two-dimensional
transverse coordinate lies in a disk of radius \(C\varepsilon\).  Coarea in
the scalar/angular witness and the remaining transverse root coordinate gives
\(C\varepsilon^{1-\alpha/4}\), hence the stated weaker uniform power.

For the complement, split relative speed at \(\varepsilon^{\alpha/8}\) and
grazing cosine at the same scale.  Gaussian velocity integration and the
collision flux give positive powers of \(\varepsilon\).  Each analytic
Jacobian minor is not identically zero on its chronological cell: moving the
chosen exclusive root or normal in a free-flight example changes the relative
position.  The one-dimensional sublevel estimate for a nonzero analytic
function, followed by Fubini over the finite atlas, gives another positive
power.  Simultaneous contacts are treated by two independent contact equations
and have higher codimension.  Taking the minimum exponent proves the lemma.
The standard labelled-tree enumeration absorbs the finite atlas and ancestry
choices.
\end{proof}

The remaining cycle contacts are controlled collectively, not assigned
fictitious independent time variables.  For a forest parameter point, let
\(N_{\rm cyc}\) be the number of later contacts between already connected
pseudo-trajectories before time \(T\).

\begin{lemma}[Ancestral contact ledger]
\label{lem:r4-b2-contact-ledger}
On every regular first-cycle witness chart and for \(|u|\le u_0\),
\[
 \int e^{uN_{\rm cyc}}
 \,d\mathfrak m_{\mathcal F}^{\rm residual}
 \le \exp\{C T e^{u}\}
\]
after the witness variable has been integrated.  The same estimate holds on
the degenerate strata after their \(O(\varepsilon^\alpha)\) activity factor is
removed.  The constant is independent of the label number and of the precise
ancestral graph.
\end{lemma}

\begin{proof}
Cut the forest-parameter space into chronological cells on which all collision
and near-collision times are ordered and every velocity is affine in the
incoming variables.  Between consecutive cell boundaries the relative
position of a fixed pair is affine in time, hence that pair has at most one
incoming contact.  Explore cells by the dependency DAG rather than by adding
formal collision variables.  When a cell boundary is crossed, either a new
label is created, which is handled by the ordinary collision-coarea operator,
or an already connected pair meets.  In the latter case the contact is a zero
of one of the transverse functions in the witness atlas.

The first such zero has already paid the power of \(\varepsilon\).  Forget it
and identify the two adjacent chronological cells.  The reflection map has
unit Jacobian in velocity and the forgetful map has uniformly bounded
multiplicity after the ancestral pair and witness type are recorded.  Repeating
this operation removes all later cycle faces without introducing a new
integration variable.  At the level of absolute activities the recursion is
\[
 Z_{m+1}(T)
 \le C\int_0^T Z_m(s)Z_0(T-s)ds,
 \qquad Z_0(T)\le e^{CT},
\]
where \(m\) counts recorded later contacts and the integral is the location of
the chronological cell boundary, not an assumed independent collision time.
The bound follows from disintegration of the already existing creation-time
coordinates.  Exponential generating functions give
\[
 \sum_{m\ge0}\frac{e^{um}}{m!}Z_m(T)
 \le\exp\{CTe^u\}.
\]
The factorial here is the combinatorial exponential-generating normalization
of the ledger recursion; it is not the volume of a nonexistent free simplex.
The same forgetful maps apply on the degenerate strata after their initial
sublevel estimate.
\end{proof}

Let a real trajectory cluster \(\lambda\) contain all labelled trajectories
and every actual contact.  For sources satisfying
\[
 |h(z([0,T]))|
 \le r(1+\sup_{t\le T}|v(t)|^2),
 \qquad
 \|\psi\|_\infty+\|\psi\|_{\rm Lip,w}\le r,
\]
put
\[
 \mathcal W_{h,\psi}(\lambda)
 =\exp\left\{
 \sum_{i\in\lambda}h(z_i([0,T]))+
 \sum_{c\in\mathcal C(\lambda)}\psi(c)
 \right\}.
\]

\begin{theorem}[All-contact forest domination]
\label{lem:r3-b2-generating}
There are \(T_0,r_0,C,\alpha>0\) such that, for \(T\le T_0\),
\(r\le r_0\), and every \(k\ge1\),
\[
 \int|d\mathfrak m_k^\varepsilon(\lambda)|
 |\mathcal W_{h,\psi}(\lambda)-1|_{\rm conn}
 \le k!C^k(T+\varepsilon)^{k-1}.
\]
For collision-source derivatives,
\[
 \sum_{j\ge0}\frac{u^j}{j!}
 |D_\psi^j\mathfrak C_k^\varepsilon(h,\psi)|
 \le k!C^k(T+\varepsilon)^{k-1}
 \exp\{CTe^{r+u}\}.
\]
The entire cyclic sector, with all its later recollisions and source marks, is
bounded by the right-hand side multiplied by
\(C\varepsilon^\alpha\).
\end{theorem}

\begin{proof}
Creation forests are treated by collision coarea: each non-root label supplies
one genuine position variable, and changing it to creation time, incoming
velocity, and normal has Jacobian
\(((v-v_*)\cdot\omega)_+\).  Gaussian energy bounds and chronological tree
summation give the forest majorant
\(k!C^k(T+\varepsilon)^{k-1}\).

For a cyclic cluster select its first cycle in chronological order and apply
Lemma~\ref{lem:r3-b2-first-cycle}.  Conditional on the resulting witness
chart, Lemma~\ref{lem:r4-b2-contact-ledger} sums every later actual contact,
including all possible collision-source selections.  The forgetful surgery
maps the residual configuration to a creation forest with uniformly bounded
multiplicity and preserves the Gaussian activity because elastic reflection
has unit velocity Jacobian and conserves energy.  Multiplication of the
first-cycle power, the residual exponential ledger, and the forest majorant
proves the cyclic estimate.  Summing the acyclic and cyclic sectors proves the
first assertion; differentiating the exact microscopic contact weight selects
contacts already counted by the ledger and proves the derivative bound.
\end{proof}

\begin{corollary}[Source-uniform all-contact cluster theorem]
\label{thm:r3-b2-gc-cluster}
The connected grand-canonical real-trajectory expansion is absolutely and
normally convergent on the declared complex source ball, with all fixed source
derivatives.  Every collision derivative is the factorial moment of the
actual incoming contact measure, and recollision-containing terms vanish
uniformly at rate \(O(\varepsilon^\alpha)\).
\end{corollary}

\begin{proof}
The logarithm of the partition function is the sum of connected labelled
clusters.  Theorem~\ref{lem:r3-b2-generating} is a geometric majorant in the
label number for sufficiently small \(T\), and Cauchy's formula justifies all
source derivatives.  On creation forests, collision coarea converges to the
Boltzmann tree integral.  The cyclic estimate removes every non-tree contact
sector uniformly.  Since the source was placed on the microscopic actual
contact list before expansion, its derivatives retain exactly those contacts.
\end{proof}
